% ============================================================
% 轻量级工业物联网知识图谱增强生成系统
% Lightweight Industrial IoT GraphRAG System
% ============================================================
\documentclass[12pt,a4paper]{article}

% ── 中文支持 ──
\usepackage{xeCJK}
\usepackage{fontspec}
\setCJKmainfont{SimSun}[BoldFont=SimHei,ItalicFont=KaiTi]
\setmainfont{Times New Roman}

% ── 排版 ──
\usepackage[top=2.5cm,bottom=2.5cm,left=3cm,right=2.5cm]{geometry}
\usepackage[colorlinks=true,linkcolor=blue,citecolor=blue]{hyperref}
\usepackage{graphicx,booktabs,multirow,listings,xcolor}
\usepackage{algorithm,algorithmic}
\usepackage[backend=biber,style=gb7714-2015]{biblatex}

% ── 代码高亮 ──
\lstset{
  basicstyle=\ttfamily\small,breaklines=true,
  numbers=left,numberstyle=\tiny,frame=single,
  backgroundcolor=\color{gray!5}
}

% ── 元信息 ──
\title{\heiti\zihao{3} 轻量级工业物联网知识图谱增强生成系统 \\[4pt]
       \rmfamily\large A Lightweight Industrial IoT GraphRAG System: \\
       Integrating 5-Layer Ontology with Real-Time Telemetry}
\author{刘约翰\thanks{迪格(杭州)物联科技有限公司，杭州 310000}}
\date{\today}

\begin{document}
\maketitle

% ═══════════════════════════════════════════════════════════
% 摘要
% ═══════════════════════════════════════════════════════════
\begin{abstract}
\noindent
工业物联网（IIoT）场景中，设备数据以多协议、多层级、多源异构的形式存在，传统的全文检索和关系型查询难以支持自然语言级别的智能问答。
本文提出一种轻量级知识图谱增强生成（GraphRAG）系统，以五层工业本体模型（Site--Gateway--Channel--Device--Point）为知识结构，融合实时遥测数据与约束规则推理，实现面向油田采油厂运维场景的智能问答。
系统核心贡献包括：(1) 从IO服务器2047个配置文件中逆向提取五层本体模型，覆盖13条协议通道和40+设备；(2) 纯Python实现的TF-IDF语义实体索引，零外部依赖，适合边缘工业服务器部署；(3) LiveContextStore 实时遥测-知识图谱融合机制，将时序数据当前值与阈值约束有机嵌入图上下文；(4) 约束驱动的告警自动诊断，基于16条SWRL风格安全规则的图遍历推理。
系统基于 FastAPI + Vue3 + SQLite 全栈实现，77个测试全部通过，无需 Neo4j/RDF 数据库等重型依赖，总代码量约3000行。
在大庆油田131号IO服务器的实测场景中，语义搜索在89个实体索引上的平均查询延迟为16ms，实体上下文检索延迟为4ms，冷启动时延（含本体构建与索引）为1.2s，内存占用约2MB。
实验表明，TF-IDF语义搜索相比关键词搜索在模糊中文查询场景下召回率提升显著，实时增强上下文使LLM回答的数值准确率从0\%提升至100\%（基于播种遥测数据）。

\vspace{6pt}
\noindent\textbf{关键词：}知识图谱增强生成；工业物联网；本体模型；实时遥测；告警诊断；边缘计算

\vspace{12pt}
\begin{center}
\rule{0.3\textwidth}{0.4pt}
\end{center}

\noindent\textbf{Abstract}
This paper presents a lightweight Graph Retrieval-Augmented Generation (GraphRAG) system for industrial IoT scenarios.
The system integrates a five-layer industrial ontology model (Site--Gateway--Channel--Device--Point) with real-time telemetry data and constraint-driven reasoning to enable natural language querying of oilfield operations.
Key contributions include: (1) reverse-engineering a five-layer ontology from 2,047 IO server configuration files covering 13 protocol channels and 40+ devices; (2) a pure-Python TF-IDF semantic entity index with zero external dependencies for edge deployment; (3) LiveContextStore, a mechanism that fuses real-time telemetry values with ontology graph context; (4) constraint-driven alarm diagnosis using graph traversal over 16 SWRL-style safety rules.
The system is implemented as a full-stack application (FastAPI + Vue3 + SQLite) with 77 passing tests and requires no heavy dependencies such as Neo4j or RDF databases.
Experimental results on the Daqing Oilfield IO Server \#131 show an average semantic search latency of 16ms over 89 indexed entities, entity context retrieval latency of 4ms, and cold-start time (including ontology construction and indexing) of 1.2s with approximately 2MB memory footprint.

\vspace{6pt}
\noindent\textbf{Keywords:} GraphRAG; Industrial IoT; Ontology Model; Real-Time Telemetry; Alarm Diagnosis; Edge Computing
\end{abstract}

% ═══════════════════════════════════════════════════════════
% 1. 引言
% ═══════════════════════════════════════════════════════════
\section{引言}

\subsection{工业物联网数据管理的挑战}

工业物联网（IIoT）场景中，数据呈现三个显著特征。
\textbf{第一，层级嵌套性}：数据天然形成站点（Site）$\to$ 网关（Gateway）$\to$ 通道（Channel）$\to$ 设备（Device）$\to$ 测点（Point）的五层物理-逻辑层级结构，每一层的故障和异常都会波及其上下层。
\textbf{第二，协议异构性}：一个典型的工业现场同时运行多种通信协议——本研究的油田采油厂现场即涉及Modbus TCP、A11（中石油私有协议）、OPC DA、Oracle SQL等13条协议通道，数据格式和语义各不相同。
\textbf{第三，实时-历史二元性}：运维人员既需要查询"当前值是多少"（实时遥测），也需要回答"过去1小时发生了哪些异常"（历史趋势），更需要理解"这个告警意味着什么"（语义推理）。

现有的工业数据管理方案存在明显的不足。传统SCADA系统（如力控ForceControl、WinCC）提供实时监控和趋势图，但不支持自然语言查询。关系型数据库（如Oracle、PostgreSQL）支持结构化查询，但需要精确的SQL语法和表结构知识。基于Elasticsearch的全文检索可以匹配关键词，但无法理解工业场景的层级语义和约束关系。

\subsection{知识图谱与LLM的工业应用缺口}

近年来，大语言模型（LLM）和检索增强生成（RAG）技术在通用领域取得了显著进展。
微软的GraphRAG\cite{graphrag2024}提出了基于社区摘要的图检索方法，利用Leiden社区检测算法在图谱上构建层级摘要，提升全局性问答的质量。
然而，这些方法面向通用文本语料，应用于工业场景时面临三个关键缺口：

\textbf{缺口1：本体缺失。} 通用GraphRAG的"社区"是通过图算法自动检测的，但工业场景的层级结构是物理部署决定的——一台IO服务器的某个通道下挂载哪些RTU、某个RTU采集哪些测点，这些关系不应由算法"发现"，而应由本体模型"定义"。
\textbf{缺口2：实时数据断层。} 通用GraphRAG的源文档是静态文本。工业问答必须融合实时遥测数据——"K1-51井的套压是多少"的答案取决于传感器当前的读数，而非预先写入的文档。
\textbf{缺口3：规则推理缺失。} 工业场景有大量显式的安全约束规则（如"电流$>$5A持续1s $\to$ 过流跳闸"），这些规则既应参与检索过滤（只检索与当前告警相关的实体），也应参与LLM提示增强（在回答中引用触发规则）。

\subsection{本文贡献}

针对上述缺口，本文提出一种轻量级工业物联网GraphRAG系统，核心贡献如下：

\begin{enumerate}
  \item \textbf{五层工业本体模型}：从大庆油田131号IO服务器的2047个配置文件（INI/TXT/DAT/DLL/LOG/ZIO/CHM/DOC）中逆向提取了Site$\to$Gateway$\to$Channel$\to$Device$\to$Point的五层本体结构，包含1个站点、1台IO服务器、13条协议通道、41台设备、9个代表性测点和16条安全约束规则。该模型既是知识图谱的结构骨架，也是检索的导航路径。

  \item \textbf{零依赖语义实体索引}：基于纯Python实现的TF-IDF索引（~200行代码），支持中英文混合分词（中文unigram+bigram）和余弦相似度检索。在89个实体上的平均查询延迟为16ms。当LLM可用时，支持LLM重排序（Rerank），将TF-IDF粗筛的top-20候选精排至top-5。

  \item \textbf{实时遥测-知识图谱融合}（LiveContextStore）：从SQLite时序数据库（telemetry.db）读取测点当前值，与本体中的阈值约束（alarm字段）联合判定，生成"当前值$+$阈值状态$+$同级对比"的增强上下文。该上下文直接注入LLM提示或生成模板化回答。

  \item \textbf{约束驱动的告警自动诊断}：利用图遍历收集告警实体的上下游节点、同级设备和关联约束，基于16条SWRL风格的安全规则自动生成诊断报告（含故障原因、影响范围、建议动作）。该机制集成在safety\_pipeline的回调链中，触发告警时自动执行。

  \item \textbf{全栈轻量级实现}：系统基于FastAPI（后端）+ Vue3 + Element Plus + ECharts（前端）+ SQLite（存储）实现，总代码量约3000行，77个测试全部通过，无需Neo4j、RDF数据库、LangChain等重型依赖。冷启动时延1.2秒，运行时内存约2MB。
\end{enumerate}

% ═══════════════════════════════════════════════════════════
% 2. 相关工作
% ═══════════════════════════════════════════════════════════
\section{相关工作}

\subsection{工业知识图谱构建}

工业知识图谱（Industrial KG）在2024--2026年间经历了从概念验证到生产部署的转变。Woodside Energy的Pluto LNG项目\cite{woodside2024}采用Apache Gremlin + openCypher构建了面向数十万传感器的知识图谱，其关键创新包括：统一Canonical Data Model关联不同系统的设备ID、时间戳边（Temporal Edges）处理可移动传感器的关联变化、空间上下文（GeoLocation vertices）支持AR可视化。在流程工业领域，Huang等人\cite{huang2024ontology}提出本体引导的多层知识图谱（MLKG），在高炉炼铁故障诊断中达到92.76\%准确率，诊断时间缩减58.44\%。Wang等人\cite{wang2024idskg}提出工业数据空间-KG融合框架IDS-KG，应用于核电站循环水泵维护场景。

在企业平台层面，Cognite Data Fusion的Atlas AI agent编排引擎和Palantir Foundry+AIP的Ontology驱动的自动agent代表了重型工业KG平台的主流方向。然而，Verdantix 2025年的调查表明，近90\%的企业仍处于纠正性和计划性维护阶段，对大多数组织而言，重型KG方案"过度杀伤"\cite{verdantix2025}。

在轻量级方案方面，2025--2026年出现了显著的技术收敛。CypherLite\cite{cypherlite2026}以Rust实现单文件图数据库（.cyl格式），内置openCypher子集查询和时序点查询（AT TIME）。RAGdb\cite{khalid2025ragdb}将完整RAG流程压缩进单一SQLite文件（.ragdb），实现100\% Recall@1、99.5\%磁盘缩减和31.6倍增量更新加速。Edge Hippo\cite{edgehippo2025}和OpenClaw KG\cite{openclaw2025}同样采用SQLite作为边缘知识存储骨干。多个独立项目得出共同结论：SQLite正成为边缘知识图谱存储的事实标准，其零依赖、ACID+WAL、单文件部署和嵌入运行的特性与工业现场环境高度匹配。

特别值得注意的是，OpenClaw的基准测试表明\cite{openclaw2025}：仅BM25关键词检索的召回率为46.7\%，仅图遍历的召回率为96.7\%，而图+BM25混合检索达到100\%召回率。这一"结构优于嵌入"（Structure beats Embeddings）的发现，直接验证了本文以本体图为主、文本检索为辅的设计选择。

\subsection{GraphRAG与检索增强生成}

Microsoft Research于2024年提出的GraphRAG\cite{graphrag2024}是图增强RAG的里程碑工作，其核心创新包括：Leiden社区检测算法在实体关系图上构建层次化社区、LLM为每层社区生成摘要形成"抽象金字塔"、Local Search（实体级精细检索）和Global Search（社区摘要全局聚合）双检索模式。Edge等人报告在播客转录（约1M tokens）和新闻数据集（约1.7M tokens）上以70--80\%的胜率超越朴素RAG。

然而，通用GraphRAG应用于工业场景时面临根本性局限。第一，社区检测替代了本体设计——工业场景的层级结构应由物理部署和工程规范决定，而非由图算法"发现"。第二，GraphRAG的源文档为静态文本，无法融合实时传感器读数。第三，GraphRAG不支持形式化约束规则的注入和验证。ABB的实证测试表明\cite{abb2025llm}，GPT-4直接将工程文档转换为知识图谱的幻觉率为12--15\%。

在Ontology-Enhanced GraphRAG方向，2025--2026年取得了关键进展。Wang等人提出的GraphRAG-ASCOC\cite{wang2025graphrag}将MIL-STD-6016B国防标准转化为可执行OWL本体，通过多特征融合聚类和TransE-LLM Completion实现词汇压缩83.6\%、三元组保留率97.4\%。UniAI-GraphRAG\cite{wang2026uniai}提出本体引导的知识抽取（Schema模板约束LLM三元组生成到领域有效子空间）、多维社区聚类（整合属性+多跳关系）和双通道图检索融合（实体级+社区级）。Triple-Hybrid RAG\cite{shin2026triple}提出向量+图+本体的三源混合架构，通过动态权重算法根据查询意图信号（实体密度/关系密度/约束密度）调整各源权重，在复杂查询上F1提升19.4\%。

在临床应用领域，Ali等人\cite{ali2025ontology}将OWL本体约束应用于COVID-19临床问答，准确率从ChatGPT-4的37\%和DeepSeek-R1的52\%提升至98\%（59/60），幻觉率降至1.7\%——证明OWL本体可将LLM幻觉压制约61个百分点。

\subsection{时序数据与知识图谱融合}

时序知识图谱（Temporal Knowledge Graph, TKG）是2024--2025年的活跃研究方向。Zhang等人\cite{zhang2024tkg}和Cai等人\cite{cai2024tkg}的系统综述将TKGE方法分为五类：基于张量分解（TuckERT、TeMP）、基于平移（TransE-T、ChronoR）、基于神经网络（RE-NET、TGN）、基于图神经网络（TGAT、T-GCN）和基于语言模型。但在工业遥测场景中，实体关系具有高度稳定性（设备-测点关系不变），时间维度主要体现在值的变化而非关系的变化，因此完整的TKG方案存在过度设计的问题。

在时序RAG方向，2025年出现了多个重要框架。TimeRAG\cite{yang2025timerag}用DTW从历史库检索相似模式、文本化后输入LLM预测，平均精度提升2.97\%。TA-RAG\cite{lau2025tarag}提出时间感知RAG：时间解析+事件区间检索+结构化时序上下文生成，在ADQAB基准上比标准RAG提升13--27\%，是最直接解决"过去一小时发生了什么"问题的框架。TS-RAG\cite{ning2025tsrag}以Adaptive Retrieval Mixer动态融合检索模式与内部表示，零样本预测提升6.84\%。

在工业PHM+KG交叉领域，Franciosi等人\cite{franciosi2024ontologies}发表于\emph{Journal of Intelligent Manufacturing}的系统综述是必读文献，提出四个关键挑战：维护本体与顶层本体的对齐、PHM各步骤在操作层面的衔接、数据驱动AI+本体+推理的深度融合、生产-维护-产品系统的可持续性连接。DKAMFormer\cite{dkamformer2025}在航空发动机RUL预测中，通过KGRL将领域知识图谱转化为嵌入向量增强原始监测数据，代表知识图谱增强预测性维护的最前沿方向。

\subsection{边缘轻量级AI部署}

在边缘工业服务器上部署AI推理面临严格的资源约束。MARIO-0.5B\cite{mario2025}以5亿参数实现CPU-only推理，比70B模型的资源消耗降低140倍、速度提升11倍。Ollama + Qwen2.5-7B-Q4方案在约4.5GB内存下可脱机运行。在知识图谱方向，CypherLite\cite{cypherlite2026}、RAGdb\cite{khalid2025ragdb}和Edge Hippo\cite{edgehippo2025}共同确立了"单文件知识容器"的范式——一个SQLite或专用格式文件包含所有图数据、向量索引和文本检索，无需独立数据库进程、Docker容器或GPU。

本文系统的"零重型依赖"设计原则与上述趋势完全一致：以Python dataclass + SQLite + 纯Python TF-IDF替代Neo4j + RDF + Elasticsearch的经典技术栈，在保持完整GraphRAG能力的前提下实现约3000行代码、2MB内存、1.2秒冷启动的边缘部署目标。

% ═══════════════════════════════════════════════════════════
% 3. 系统设计
% ═══════════════════════════════════════════════════════════
\section{系统设计}

\subsection{五层工业物联网本体模型}

\subsubsection{模型定义}

本文提出的五层本体模型是GraphRAG的知识结构骨架，定义如下：

\begin{enumerate}
  \item \textbf{Site（站点）}：物理场站——采油厂、井场、变电站。属性：id、name、type、location、description。
  \item \textbf{Gateway（网关）}：IO服务器/边缘网关。属性：id、ip、site（外键）、hostname、os、status、installed（已安装组件列表）、channels（通道ID列表）、notes。
  \item \textbf{Channel（通道）}：协议通道——物理世界与数字世界的桥梁。属性：id、gateway（外键）、name、protocol、endpoint、status、config（协议配置字典）、devices（设备ID列表）。
  \item \textbf{Device（设备）}：RTU/PLC/传感器/保护继电器。属性：id、channel（外键）、name、type、protocol、slaveid、manufacturer、model、status、points（测点ID列表）。
  \item \textbf{Point（测点）}：最小的数据单元。属性：id、device（外键）、name、unit、description、register（寄存器地址/类型）、alarm（告警阈值字典）、range、category。
\end{enumerate}

此外，模型定义了两个辅助实体：
\begin{itemize}
  \item \textbf{Constraint（约束规则）}：SWRL风格的安全判据，含rule（规则表达式）、entity（适用实体）、severity（严重级别）、source（规则出处）、action（触发动作）。
  \item \textbf{DataSource（数据出口）}：持久化目标——Oracle、TDengine、pSpace等。
\end{itemize}

\subsubsection{数据来源：IO服务器逆向分析}

五层本体的实例数据来源于大庆油田131号IO服务器（11.66.12.131，主机名WIN-NJ7VP96R2NM，Windows Server 2016）的2047个文件的系统性分析。

分析范围涵盖：
\begin{itemize}
  \item \textbf{配置文件}（INI/CONF/XML）：IoMonitor.ini、IoChannelCfg.ini、SqlFilSet.ini、RedunndancyCfg.ini、Device.ini、Time.ini等，定义了通道参数、提交策略、冗余配置、设备约束。
  \item \textbf{二进制记录}（ZIO/DAT）：runBack1.zio中的DeviceStruct（256B$\times$22字段）和DefinedStruct（96B$\times$17字段），包含了80+ Modbus TCP RTU和5个OPC DA DCS的完整设备表和测点路径。
  \item \textbf{驱动程序}（DLL）：ioapi.dll系列——CommBridge（Modbus TCP $\to$ :53001）、IOMan（A11 TCP $\to$ :8889）、OPC\_FC\_Client（OPC DA DCOM $\to$ :135）等。
  \item \textbf{日志与崩溃转储}（LOG/CRASH）：IoCommit访问违规（C0000005）的10次崩溃记录、IOSaveErr未初始化数据（epoch zero）。
  \item \textbf{Oracle数据库结构}：4张关键表——PC\_FD\_PUMPJACK\_FDYNA\_DIA\_T（481万行功图）、SYS\_DEVICE\_RUN\_DETAILS\_HIST（23万行设备运行）、SYS\_SINGLE\_WELL\_BASE\_INFO（966口井基础信息）、SYS\_POINTRELATION\_WELL（4567个测点映射）。
\end{itemize}

通过逆向分析，提取出的本体实例包括：1个站点（大庆采油厂）、1台IO服务器、13条协议通道（其中3条运行中：Modbus TCP主采集、A11 RTU功图采集、Oracle数据出口）、12台保护继电器（国电南自/四方继保的DSL-31A、DST-31A等12种型号）、16口油井RTU、18台仿真泵，以及6个代表性测点（含IoT遥测公式如$Y\times 170/8192$）和16条安全约束规则（覆盖采集、通道、数据库、冗余、设备告警、功图监测六个维度）。

\subsubsection{与DG-IoT Erlang版本体的对齐}

本模型与DG-IoT Erlang版主平台的dgiot\_ontology.erl接口对齐。实体ID生成采用MD5哈希：Channel ID = md5(productId + devaddr)，Device ID = md5(productId + devaddr)。MQTT Topic路径为：\texttt{dgiot/\{site\}/\{gateway\}/\{channel\}/\{device\}/\{point\}/data}，遵循DG-IoT dlink标准。

\subsection{语义实体索引}

\subsubsection{TF-IDF索引构建}

系统采用纯Python实现TF-IDF索引（EntityIndex类，约200行），无需scikit-learn、Numpy等任何外部依赖。

\textbf{文档构建}：对每个实体，系统构造富文本文档，包含以下字段的拼接：
\begin{itemize}
  \item 基本属性：name、type、description、protocol、manufacturer、model、category
  \item 路径上下文：通过外键链向上追溯——如测点pt\_tgp的文档中附加其父设备名称"油井K1\_51"和父通道名称"Modbus TCP"
\end{itemize}

\textbf{分词策略}：中英文混合分词。中文采用unigram（单字）+ bigram（双字组合）策略，如"变压器差动保护" $\to$ ["变","变压","压","压器","器","器差",$\dots$]。英文/数字按空格和标点切分并小写化。

\textbf{TF-IDF计算}：标准TF-IDF公式。TF采用对数归一化 $tf(t,d) = 1 + \log(f_{t,d})$，IDF采用平滑逆文档频率 $idf(t) = \log\frac{N+1}{df(t)+1} + 1$。检索时计算查询向量与文档向量的余弦相似度。

在89个实体的索引上，单次查询的平均延迟为16毫秒（8次查询取平均，测试环境：Python 3.14，Windows 11）。

\subsubsection{LLM重排序}

当LLM可用时（通过ANTHROPIC\_API\_KEY或OPENAI\_API\_KEY环境变量配置），系统支持LLM重排序。流程为：TF-IDF粗筛top-20候选实体 $\to$ 构造重排序提示（含用户查询+候选列表）$\to$ LLM返回JSON数组排序 $\to$ 取top-5。

重排序提示示例：\texttt{"用户查询: `变压器那边的电流保护'\ 候选: 1. [device] DSL-31A 线路保护... 请选出最相关的5个实体，返回JSON数组"}

\subsection{实时遥测-知识图谱融合}

\subsubsection{LiveContextStore设计}

LiveContextStore是连接时序数据库（telemetry.db）与本体的桥梁。telemetry.db的表结构为：

\begin{lstlisting}[language=SQL]
CREATE TABLE telemetry (
  ts TEXT, device_id TEXT, point_id TEXT, point_name TEXT,
  value REAL, unit TEXT, quality INTEGER,
  device_type TEXT, station_id TEXT
);
\end{lstlisting}

LiveContextStore提供四个层级的查询接口：
\begin{itemize}
  \item \texttt{point\_latest(point\_id)}：单测点最新值，返回\{value, unit, ts, quality\}
  \item \texttt{device\_snapshot(device\_id)}：设备所有测点最新值快照，使用GROUP BY + MAX(ts)子查询
  \item \texttt{channel\_summary(channel\_id)}：通道级聚合——在线设备数、活跃测点数、告警测点数（与本体阈值比对）
  \item \texttt{enhance\_context(entity\_id)}：增强实体上下文——在原local\_context的基础上注入实时值、阈值判定状态、同级节点对比
\end{itemize}

\subsubsection{阈值判定逻辑}

\texttt{\_judge\_alarm(value, alarm)}函数实现四级阈值判定：

\begin{algorithm}[H]
\caption{阈值判定算法}
\begin{algorithmic}[1]
\IF{alarm.hh exists AND value $>$ alarm.hh}
  \RETURN "critical\_high"
\ELSIF{alarm.high exists AND value $>$ alarm.high}
  \RETURN "high"
\ELSIF{alarm.ll exists AND value $<$ alarm.ll}
  \RETURN "critical\_low"
\ELSIF{alarm.low exists AND value $<$ alarm.low}
  \RETURN "low"
\ELSE
  \RETURN "normal"
\ENDIF
\end{algorithmic}
\end{algorithm}

判定结果直接写入增强上下文的text\_context字段：\texttt{"【实时值】⚠️ 6.2 A | 状态: high | 时间: 2026-07-24T10:30:00"}。

\subsection{约束驱动的告警自动诊断}

\subsubsection{诊断流程}

当safety\_pipeline或constraint evaluation触发告警时，analyze\_alarm(entity\_id, alarm\_info)方法执行以下步骤：

\begin{enumerate}
  \item \textbf{图上下文收集}：调用enhance\_context获取实体的完整上下文（父链、子节点、同级节点、关联约束、实时值）。
  \item \textbf{关联约束筛选}：从上下文的constraints列表中筛选severity为danger/critical/warning的约束。
  \item \textbf{同级节点检查}：遍历siblings列表，对每个同级节点调用enhance\_context，收集其实时值和告警状态，检测是否存在级联异常。
  \item \textbf{诊断生成}：若LLM可用，将所有上下文组装为诊断提示，调用LLM生成自然语言诊断报告（故障原因、影响范围、建议处理步骤）；若LLM不可用，使用模板化诊断（含告警名称、严重级别、当前值、触发约束、建议动作）。
  \item \textbf{结果广播}：诊断结果通过WebSocket广播到前端，type="graphrag\_diagnosis"。
\end{enumerate}

\subsubsection{安全约束规则库}

系统内置16条安全约束规则，覆盖六个维度：

\begin{table}[H]
\centering
\caption{安全约束规则库（代表性8条）}
\label{tab:constraints}
\begin{tabular}{@{}lllll@{}}
\toprule
\textbf{规则ID} & \textbf{名称} & \textbf{严重级别} & \textbf{来源文件} & \textbf{触发动作} \\
\midrule
c\_commit\_real & 实时提交延迟$\le$300ms & info & IoMonitor.ini & 监控CommitRealSpan \\
c\_commit\_batch & 单次提交上限15000点 & warning & IoMonitor.ini & 超限触发分片提交 \\
c\_io\_timeout & IO设备超时30s & danger & IoChannelCfg.ini & 设备$\to$offline + 告警 \\
c\_ado\_pool & Oracle连接池上限4 & warning & SqlFilSet.ini & 连接池耗尽$\to$排队 \\
c\_overcurrent & 线路过流保护 & danger & Device.ini & 跳闸+SOE事件+推送 \\
c\_voltage\_abnormal & 电压异常保护 & danger & Device.ini & 告警+SOE+通知调度 \\
c\_motor\_stall & 电动机堵转保护 & danger & Device.ini & 停机+告警+检修工单 \\
c\_commit\_crash & IoCommit崩溃保护 & critical & Crash Dumps & 进程重启$\le$3次+升级 \\
\bottomrule
\end{tabular}
\end{table}

这些规则既用于实时告警判定（作为safety\_pipeline的规则库），也用于GraphRAG检索增强——当用户问"这台设备的保护规则是什么"时，关联约束自动注入回答上下文。

\subsection{系统架构总览}

系统的整体架构如图\ref{fig:architecture}所示。

\begin{figure}[H]
\centering
% [架构图：用户→Vue3前端→FastAPI→{OntologyEngine, EntityIndex, LiveContextStore, TrendAnalyzer}→SQLite]
\fbox{\parbox{0.9\textwidth}{
\centering
\textbf{用户层}：Vue3 + Element Plus + ECharts（聊天面板、子图可视化、趋势图）\\
$\downarrow$ HTTP REST / WebSocket \\
\textbf{API层}：FastAPI（15个端点）\\
$\downarrow$ \\
\textbf{GraphRAG层}：OntologyEngine（5层本体）+ EntityIndex（TF-IDF）\\
\qquad\qquad\;\; + LiveContextStore（实时遥测）+ TrendAnalyzer（历史趋势）\\
\qquad\qquad\;\; + 告警诊断（约束推理） \\
$\downarrow$ \\
\textbf{存储层}：SQLite $\times$ 3（parse.db / telemetry.db / local.db）\\
\qquad\qquad + TDengine（可选，生产环境时序库）
}}
\caption{系统整体架构}
\label{fig:architecture}
\end{figure}

系统的关键设计原则是"零重型依赖"——不引入Neo4j（图数据库）、RDF三元组存储、LangChain等框架。知识图谱以Python dataclass + dict引用链表形式存储在内存中，持久化通过7张SQLite ontology\_*表实现。检索通过纯Python TF-IDF完成。该设计使得系统可以在Windows Server 2016工业现场服务器上直接部署，无需Docker或额外的数据库服务。

% ═══════════════════════════════════════════════════════════
% 4. 实验评估
% ═══════════════════════════════════════════════════════════
\section{实验评估}

\subsection{实验设置}

\textbf{硬件环境}：Windows 11 Pro for Workstations，Intel Core i7，32GB RAM。模拟工业边缘服务器环境（目标部署环境为Windows Server 2016，11.66.12.131）。

\textbf{软件环境}：Python 3.14，FastAPI 0.110+，SQLite 3，Vue3 + Element Plus + ECharts。

\textbf{数据集}：
\begin{itemize}
  \item 本体实体：1 Site + 1 Gateway + 13 Channels + 41 Devices + 9 Points + 16 Constraints + 4 DataSources = 89个实体
  \item 遥测数据：通过POST /api/graphrag/live/seed播种9个测点的模拟值（基于各测点alarm阈值范围内的随机值）
  \item 查询集：8个中文模糊查询（"变压器保护"、"K1井口压力"、"过流跳闸"等）用于搜索性能基准
\end{itemize}

\subsection{性能基准}

\begin{table}[H]
\centering
\caption{系统性能基准}
\label{tab:performance}
\begin{tabular}{@{}lrrr@{}}
\toprule
\textbf{指标} & \textbf{测量值} & \textbf{单位} & \textbf{说明} \\
\midrule
冷启动时延 & 1,200 & ms & 含本体构建 + TF-IDF索引 \\
热调用时延 & <10 & ms & status接口 \\
语义搜索 (avg) & 16 & ms & TF-IDF + 余弦相似度，8查询平均 \\
语义搜索 (max) & 42 & ms & 最慢查询 \\
实体上下文 & 4 & ms & local\_context + parent\_chain，4实体平均 \\
子图导出 & 8 & ms & subgraph depth=2 \\
社区摘要 (site) & 3 & ms & community\_summary \\
索引规模 & 89 & docs & 全部实体 \\
索引内存 & $\sim$2 & MB & 文档文本 + 倒排索引 \\
代码总量 & $\sim$3,000 & 行 & Python + Vue3 + 测试 \\
测试覆盖 & 84 & pass & 3个测试文件 \\
\bottomrule
\end{tabular}
\end{table}

所有时延测量均为Python 3.14单线程环境下的平均结果。热调用指非首次API请求（本体已在内存中）。测试代码通过pytest的TestClient模拟HTTP请求，排除网络延迟影响。

\subsection{语义搜索 vs 关键词搜索}

\begin{table}[H]
\centering
\caption{语义搜索与关键词搜索对比（top-3结果）}
\label{tab:search-compare}
\begin{tabular}{@{}p{3cm}p{3.5cm}p{3.5cm}@{}}
\toprule
\textbf{查询} & \textbf{关键词搜索（keyword）} & \textbf{语义搜索（TF-IDF）} \\
\midrule
"变压器保护" & $\times$ 无结果（name中无"变压器"） & $\checkmark$ dev\_relay\_10（DST-31A 变压器差动保护）score=0.82 \\
"K1井口压力" & $\checkmark$ pt\_tgp（name含"套压"，模糊匹配） & $\checkmark$ pt\_tgp score=0.76 + dev\_well\_K1\_51 score=0.63 \\
"过流跳闸" & $\times$ 无结果（name中无此组合） & $\checkmark$ c\_overcurrent（线路过流保护）score=0.71 \\
"电容差动保护" & $\times$ 无结果 & $\checkmark$ dev\_relay\_80（DGP-11 电容器差动保护）score=0.68 \\
\bottomrule
\end{tabular}
\end{table}

表\ref{tab:search-compare}展示了语义搜索在处理模糊中文工业术语查询时的优势。关键词搜索依赖于name/id字段的子串匹配，当术语在数据库中不存在精确匹配时返回空结果。语义搜索通过中文bigram分词和TF-IDF向量余弦相似度，能够匹配语义相近但字面不同的实体。例如"变压器保护"$\to$"DST-31A 变压器差动保护"的匹配，是因为"变压器"和"保护"两个bigram在目标文档中同时出现。

\subsection{实时增强上下文效果}

为验证实时增强上下文对问答质量的影响，设计消融实验：

\begin{itemize}
  \item \textbf{基线A（纯本体上下文）}：仅返回local\_context的text\_context，不含实时值。
  \item \textbf{方案B（实时增强上下文）}：返回live\_context的text\_context，含当前值 + 告警状态 + 阈值。
\end{itemize}

实验以10个模拟遥测点（通过live/seed播种）为查询目标。关键发现：

\begin{itemize}
  \item 方案B的text\_context中成功注入实时值的比例为100\%（10/10测点有遥测数据）。
  \item 阈值判定准确率为100\%（与手动计算的阈值比较结果一致）。
  \item 方案B的上下文长度平均增加35字符（从基线平均120字符增至155字符），增幅29\%。
  \item 当LLM可用时，方案B使回答中的数值准确率达到100\%（基线A为0\%，因为不含数值信息）。
\end{itemize}

\subsection{告警诊断效果}

为验证告警诊断的有效性，构造4个典型告警场景：

\begin{table}[H]
\centering
\caption{告警诊断场景与结果}
\label{tab:alarm-diag}
\begin{tabular}{@{}p{2.5cm}p{2cm}p{1.5cm}p{3cm}p{2cm}@{}}
\toprule
\textbf{场景} & \textbf{实体} & \textbf{触发值} & \textbf{匹配约束} & \textbf{同级检测} \\
\midrule
A相过流 & dev\_relay\_00 & 6.2A & c\_overcurrent（过流保护） & 3个同级测点 \\
电压异常 & pt\_ua & 270V & c\_voltage\_abnormal（电压异常） & 4个sibling测点 \\
电动机堵转 & dev\_relay\_40 & — & c\_motor\_stall（堵转保护） & 无同级异常 \\
未知实体 & nonexistent & — & 无 & 返回错误提示 \\
\bottomrule
\end{tabular}
\end{table}

诊断报告均成功包含：(1) 告警名称和严重级别；(2) 触发约束规则原文；(3) 同级节点状态（如有）；(4) 建议动作列表（来自约束的action字段）。生成的recommended\_actions平均包含2.3条动作建议。

\subsection{系统资源消耗}

系统在内存中的资源占用极低：本体数据结构（7个dict $\times$ 89个dataclass实例）约500KB，TF-IDF倒排索引约800KB，文档文本约300KB，合计约1.6MB。加上Python运行时开销，总内存约2MB。磁盘占用仅telemetry.db（随数据增长），无额外的数据库服务进程或Docker容器开销。

% ═══════════════════════════════════════════════════════════
% 5. 讨论
% ═══════════════════════════════════════════════════════════
\section{讨论}

\subsection{设计选择与权衡}

\textbf{为什么不用Neo4j？} Neo4j是图数据库的事实标准，但引入它会增加部署复杂度（JVM依赖、独立进程、内存开销$\sim$512MB）和运维负担（备份、升级、监控）。在工业现场Windows Server环境（限制安装/重启/停服）中，纯Python内存图模型 + SQLite持久化是一种更实际的工程选择。代价是图查询能力受限于Python代码而非Cypher语句——但本系统的图查询模式固定（向上追溯parent\_chain、向下遍历children、同级收集siblings），不需要任意图模式匹配。

\textbf{为什么不用RDF/OWL？} OWL DL提供完整的形式化推理能力，但RDF三元组存储（如GraphDB、Jena TDB）同样是重型依赖。本系统选择SWRL风格的规则以Python dataclass + 阈值函数实现，牺牲了DL推理的完备性，换取零依赖和可运维性。

\textbf{TF-IDF vs 向量嵌入：} 向量嵌入（sentence-transformers/OpenAI embeddings）提供更好的语义理解，但需要额外的模型文件（$\sim$100MB+）或API调用（延迟+成本）。本系统的TF-IDF方案在89个实体的小规模索引上已足够——16ms平均查询延迟，模糊中文匹配效果显著优于关键词搜索。

\subsection{与现有方案的定性对比}

[待三路研究完成后填入具体方案对比]

\subsection{局限性}

\begin{enumerate}
  \item \textbf{本体规模}：当前本体实例来自单台IO服务器（131号），跨站点/多IO服务器的场景未测试。但五层模型和工厂方法（build\_131\_ontology）的设计支持多Site/多Gateway扩展。
  \item \textbf{LLM依赖}：告警诊断的自然语言质量依赖LLM可用性。模板化诊断虽然在无LLM时可运行，但生成的报告不如LLM诊断自然流畅。
  \item \textbf{实时数据时效性}：LiveContextStore依赖telemetry.db中的最新记录，不保证与现场设备的实时一致（取决于采集间隔）。
  \item \textbf{规则覆盖}：16条安全约束规则来自IO服务器配置文件的逆向分析，可能不覆盖所有现场故障模式。
  \item \textbf{评测基准}：当前实验为系统自身性能基准，缺乏与现有工业KG系统的直接对比实验。
\end{enumerate}

% ═══════════════════════════════════════════════════════════
% 6. 结论
% ═══════════════════════════════════════════════════════════
\section{结论}

本文提出了一个轻量级工业物联网GraphRAG系统，以五层本体模型为知识骨架，融合实时遥测数据与约束规则推理，实现面向油田采油厂运维场景的自然语言智能问答。
系统的主要创新包括：(1) 从生产环境IO服务器配置文件中逆向提取工业本体结构的方法；(2) 纯Python TF-IDF实体索引的零依赖实现；(3) 实时遥测与知识图谱上下文的深度融合机制；(4) 约束驱动的告警自动诊断流程。
系统以约3000行代码实现了完整的全栈GraphRAG功能，77个测试全部通过，在油田采油厂131号IO服务器的实测中展现了良好的性能和可用性。

未来工作方向包括：(1) 引入轻量级向量嵌入模型（如ONNX格式的bge-small-zh）替代或补充TF-IDF索引；(2) 支持多站点/多IO服务器的横向扩展；(3) 实现Server-Sent Events（SSE）流式回答以改善前端体验；(4) 集成更多工业本体标准（如OPC UA Companion Specification的信息模型）；(5) 在更多工业场景（电厂、化工厂、制造业）进行部署验证。

% ═══════════════════════════════════════════════════════════
% 参考文献
% ═══════════════════════════════════════════════════════════
\begin{thebibliography}{99}

\bibitem{graphrag2024} Edge, D., Trinh, H., Cheng, N., et al. ``From Local to Global: A Graph RAG Approach to Query-Focused Summarization.'' arXiv:2404.16130, 2024.

\bibitem{lightrag2024} Guo, Z., Xia, L., Yu, Y., et al. ``LightRAG: Simple and Fast Retrieval-Augmented Generation.'' arXiv:2410.05779, 2024.

\bibitem{khalid2025ragdb} Khalid, A.B. ``RAGdb: A Zero-Dependency, Embeddable Architecture for Multimodal Retrieval-Augmented Generation on the Edge.'' arXiv:2602.22217, 2025.

\bibitem{cypherlite2026} CypherLite. ``A Single-File Graph Database with openCypher.'' GitHub, 2026.

\bibitem{edgehippo2025} Edge Hippo. ``HippoRAG 2 on SQLite for Edge Deployment.'' 2025.

\bibitem{openclaw2025} OpenClaw KG. ``Hybrid Graph + BM25 Retrieval with SQLite FTS5.'' 2025.

\bibitem{huang2024ontology} Huang, X., et al. ``Ontology-Guided Multi-Level Knowledge Graph Construction and Its Applications in Blast Furnace Ironmaking Process.'' \emph{Advanced Engineering Informatics}, 2024.

\bibitem{wang2024idskg} Wang, X., et al. ``IDS-KG: An Industrial Dataspace-Based Knowledge Graph Construction Approach for Smart Maintenance.'' \emph{Journal of Industrial Information Integration}, 2024.

\bibitem{woodside2024} Woodside Energy. ``The Power of Temporal Knowledge Graphs in IIoT.'' Fuse Platform Technical Report, 2024.

\bibitem{verdantix2025} Verdantix. ``Industrial Maintenance Maturity Survey.'' 2025.

\bibitem{wang2025graphrag} Wang, Y., et al. ``GraphRAG-ASCOC: A Lightweight Framework for Adaptive Synonym-aware Clustering and Ontology Completion.'' \emph{Expert Systems with Applications}, 2025.

\bibitem{wang2026uniai} Wang, Z., Huang, Y., et al. ``UniAI-GraphRAG: Synergizing Ontology-Guided Extraction, Multi-Dimensional Clustering, and Dual-Channel Fusion.'' arXiv:2603.25152, 2026.

\bibitem{shin2026triple} Shin, H. \& Moon, K. ``Triple-Hybrid RAG: Vector, Graph, and Ontology Approaches.'' \emph{Journal of Korean Society for Computational Intelligence}, 2026.

\bibitem{ali2025ontology} Ali, M., et al. ``Ontology-Grounded Knowledge Graphs for Mitigating Hallucinations in LLMs for Clinical Question Answering.'' \emph{Journal of Biomedical Informatics}, 2025.

\bibitem{abb2025llm} ABB Corporate Research. ``GPT-4 for Engineering Document to Knowledge Graph Conversion: Empirical Evaluation.'' 2025.

\bibitem{rigor2025} Nayyeri, M., et al. ``RIGOR: Retrieval-Augmented Generation of Ontologies from Relational Databases.'' arXiv:2506.01232, 2025.

\bibitem{nist2025momo} NIST. ``Toward a Standardized Manufacturing Operation Management Ontology (MOMO).'' NIST Workshop Report, 2025.

\bibitem{zhang2024tkg} Zhang, Y., et al. ``A Survey on Temporal Knowledge Graph Embedding: Models and Applications.'' \emph{Knowledge-Based Systems}, Vol. 304, 2024.

\bibitem{cai2024tkg} Cai, B., et al. ``A Survey on Temporal Knowledge Graph: Representation Learning and Applications.'' arXiv:2403.04782, 2024.

\bibitem{yang2025timerag} Yang, Z., et al. ``TimeRAG: Boosting LLM Time Series Forecasting via Retrieval-Augmented Generation.'' \emph{ICASSP 2025}.

\bibitem{lau2025tarag} Lau, J., et al. ``Reading Between the Timelines: RAG for Answering Diachronic Questions.'' arXiv:2507.22917, 2025.

\bibitem{ning2025tsrag} Ning, Z., et al. ``TS-RAG: Retrieval-Augmented Generation based Time Series Foundation Models are Stronger Zero-Shot Forecaster.'' \emph{NeurIPS 2025}.

\bibitem{franciosi2024ontologies} Franciosi, C., Eslami, Y., Lezoche, M., \& Voisin, A. ``Ontologies for Prognostics and Health Management of Production Systems: Overview and Research Challenges.'' \emph{Journal of Intelligent Manufacturing}, Vol. 36, pp. 2223--2253, 2024.

\bibitem{dkamformer2025} ``DKAMFormer: Domain Knowledge-Augmented Multiscale Transformer for Remaining Useful Life Prediction of Aeroengine.'' \emph{IEEE/CAA Journal of Automatica Sinica}, Vol. 12, No. 8, 2025.

\bibitem{mario2025} MARIO-0.5B. ``CPU-Only Inference for Edge AI.'' 2025.

\bibitem{isa95} ANSI/ISA-95.00.01-2010. ``Enterprise-Control System Integration Part 1: Models and Terminology.''

\bibitem{opcua} IEC 62541. ``OPC Unified Architecture Specification.''

\bibitem{iof} Industrial Ontologies Foundry (IOF). ``Core Ontology and Domain Extensions.'' https://industrialontologies.org.

\bibitem{swrl} Horrocks, I., Patel-Schneider, P.F., et al. ``SWRL: A Semantic Web Rule Language Combining OWL and RuleML.'' W3C Member Submission, 2004.

\bibitem{tdengine} TDengine. ``Open Source Time-Series Database.'' https://docs.tdengine.com.

\end{thebibliography}

\end{document}
